\chapter{A finite-energy inversion with a nonsmooth closing point}\label{sec:examples-inverted-catenoid}
The catenoid of Example \ref{ex:explicit-catenoid} and the inversion identity of Section \ref{sec:willmore-multiplicity-inversion} give an explicit warning about compactification. The calculations use the first-variation boundary flux in \cite{SimW93}, equation (1.5), but do not assert a removable-singularity theorem.

\begin{proposition}[The inverted catenoid has energy eight pi]\label{prop:inverted-catenoid-energy}
For $a>0$, let
\[
 C_a(\theta,v)=(a\cosh v\cos\theta,a\cosh v\sin\theta,av),
 \qquad \theta\in\mathbb R/2\pi\mathbb Z,\ v\in\mathbb R.
\]
It avoids zero, and $I_0\circ C_a$ is a smooth embedded noncompact surface of finite Willmore energy $8\pi$. For its portion with $|v|<T$, $T>0$, the energy is
\[
 8\pi\frac{\cosh T\sinh T+T}{\cosh^2T+T^2}.
\]
Both ends converge to the origin in the target.
\end{proposition}
\begin{proof}
The original catenoid is embedded and minimal, and $|C_a|^2=a^2(\cosh^2v+v^2)\geq a^2$, so inversion is a diffeomorphism on its image. Lemma \ref{lem:inversion-mean-curvature-vector} expresses the inverted energy on $|v|<T$ as $4\int |C_a^\perp|^2/|C_a|^4$. Apply Theorem \ref{thm:willmore-radial-square-identity} to the compact truncated catenoid with pole zero. It has no pole preimages and zero original mean curvature, so this energy is twice its boundary flux.

At $v=T$, the outward conormal is $C_{a,v}/(a\cosh T)$ and the boundary arclength is $a\cosh T\,d\theta$. Therefore that circle contributes
\[
 \int_0^{2\pi}\frac{C_a\cdot C_{a,v}}{|C_a|^2}d\theta
 =2\pi\frac{\cosh T\sinh T+T}{\cosh^2T+T^2}.
\]
The lower boundary gives the same value, because both $v$ and the outward derivative reverse sign. Multiplying the sum by two proves the displayed energy. Increasing $T$ exhausts the parameter surface, so monotone convergence of the nonnegative energy density and the ratio limit one give $8\pi$. Finally $|I_0(C_a)|=1/|C_a|\to0$ as $v\to\pm\infty$.
\end{proof}

\begin{proposition}[The closing sheets have no twice-differentiable extension]\label{prop:inverted-catenoid-not-c2}
Each end of the inverted catenoid can be filled at zero as a rotational graph that is continuously differentiable with horizontal tangent plane, but it has no $C^2$ graph extension there. In particular its finite energy does not make the closure a smooth closed immersed surface to which the smooth Li--Yau theorem can simply be applied.
\end{proposition}
\begin{proof}
For the upper end, the radial and height coordinates are
\[
 \rho(v)=\frac{\cosh v}{a(\cosh^2v+v^2)},\qquad
 z(v)=\frac{v}{a(\cosh^2v+v^2)}.
\]
As $v\to+\infty$, they satisfy
\[
 \rho(v)=\frac2a e^{-v}(1+o(1)),\qquad
 z(v)=\frac4a v e^{-2v}(1+o(1)),
 \qquad \frac{z(v)}{\rho(v)^2}=av\left(1+\frac{v^2}{\cosh^2v}\right)\longrightarrow\infty.
\]
Differentiating the exact rational expressions shows $\rho'(v)=-(2/a)e^{-v}(1+o(1))<0$ eventually and $z'(v)/\rho'(v)=O(ve^{-v})\to0$. Thus the end is a graph $z=f(\rho)$ for small positive $\rho$, with $f(0)=0$ and radial gradient tending to zero. This gives a $C^1$ extension on the planar disk. If it had a $C^2$ extension, its zero value and gradient at the origin and Taylor's theorem would imply $f(\rho)=O(\rho^2)$, contradicting the displayed divergent quotient. The lower end is its reflection and has the same conclusion. The two closing sheets therefore cannot be treated as smooth parameter sheets at their common limiting point.
\end{proof}
